\newcommand{\figuraCampoMarea}{%
\begin{figure}[H]
    \centering
    \begin{tikzpicture}[>=Latex,font=\small,line cap=round,line join=round]
        \coordinate (A) at (3.0,3.0);
        \coordinate (B) at (9.2,3.0);

        \node[font=\bfseries] at (3.0,5.05)
            {(a) Gravedad lunar};
        \node[font=\bfseries] at (9.2,5.05)
            {(b) Campo de marea};

        \foreach \P in {A,B} {
            \fill[blue!13] (\P) circle (1.35);
            \draw[very thick,blue!65!black] (\P) circle (1.35);
            \fill[black] (\P) circle (0.05);
        }

        \draw[->,thick] (0.8,4.55)--(5.2,4.55)
            node[right] {hacia la Luna};

        \draw[->,ultra thick,orange!80!black] (1.65,3.0)--(2.22,3.0);
        \draw[->,ultra thick,orange!80!black] (3.0,3.0)--(3.92,3.0);
        \draw[->,ultra thick,orange!80!black] (4.35,3.0)--(5.72,3.0);
        \draw[->,ultra thick,orange!80!black] (3.0,4.35)--(4.02,4.14);
        \draw[->,ultra thick,orange!80!black] (3.0,1.65)--(4.02,1.86);
        \node[align=center,text=orange!80!black] at (3.0,0.92)
            {todos los puntos aceleran\\hacia la Luna};

        \draw[->,ultra thick,red!75!black] (10.55,3.0)--(11.72,3.0);
        \draw[->,ultra thick,red!75!black] (7.85,3.0)--(6.68,3.0);
        \draw[->,ultra thick,red!75!black] (9.2,4.35)--(9.2,3.68);
        \draw[->,ultra thick,red!75!black] (9.2,1.65)--(9.2,2.32);
        \node[align=center,text=red!75!black] at (9.2,0.92)
            {estiramiento axial\\y compresión transversal};

        \draw[densely dashed,black!45] (6.2,0.55)--(6.2,5.35);
        \node[align=center,text=black!70] at (6.2,5.72)
            {$\vec g_{\mathrm{tidal}}
            =\vec g_{\mathrm{Luna}}-\vec g_{\mathrm{Luna}}(O)$};
    \end{tikzpicture}
    \caption{La gravedad lunar (a) apunta hacia la Luna en toda la Tierra,
    pero su magnitud es mayor en el lado cercano. Al restar la aceleración del
    centro terrestre se obtiene el campo de marea (b): los puntos sublunar y
    antilunar se separan del centro, mientras que las regiones transversales
    son aceleradas hacia el eje Tierra--Luna. Las flechas están exageradas.}
    \label{fig:campo-marea}
\end{figure}%
}